\suppnote{Fields, road registration, estimation and uncertainty}{note:estimation}

\paragraph{Discrete fields and conservation.} Let $F_t\in\R^{n_f}$ be the observed signed crossing field on interior faces; a prediction is distinguished by a hat. The cell--face incidence matrix $D\in\R^{n_c\times n_f}$ maps face flux to net cell outflow under the frozen orientation. Thus
\begin{equation}
 x_t=DF_t,\qquad P_0=I-D^\dagger D,\qquad
 \Fcirc_t=P_0F_t,\qquad D\Fcirc_t=0 .
 \label{eq:supp-field-decomposition}
\end{equation}
The observed source--sink balance and divergence-free routing field give the spatially orthogonal decomposition $F_t=D^\dagger x_t+\Fcirc_t$. Orthogonality at an instant does not imply statistical independence across observations. The minimum-norm lift uses a gauge-pinned Laplacian solve rather than an explicitly formed Moore--Penrose inverse; on the zero-sum compatible source space the two constructions coincide. Conservation is imposed by construction or projection, not learned.

\paragraph{Road registration.} The registered road graph uses road nodes and registered road segments as edges; it is not trimmed to the mobility footprint. If $B_R$ is its node--edge incidence, $A_R$ lifts signed road-edge flows to interior faces, and $\Pi$ aggregates cells to nodes, registration satisfies $DA_R=\Pi B_R$ in integer arithmetic. For a cycle basis $Z_R$ with $B_RZ_R=0$,
\begin{equation}
 \UR=\operatorname{orth}(A_RZ_R),\qquad
 \UR^\top\UR=I,\qquad D\UR=0 .
 \label{eq:supp-road-basis}
\end{equation}
The signed entries of $A_R$ count open-segment crossings in the frozen face orientation. Geometry, study domain and road registration determine the basis; the analysis footprint mask, mobility values, dates and folds do not. This independence is conditional on the registered spatial conventions. The observed road-coordinate state is $y_t=\UR^\top\Fcirc_t$. The physical lift $\CR=\UR K_R\Px$, with $\Px=DD^\dagger$, is used for the local response tensors in Note~\ref{note:spatial}.

\paragraph{Response estimation.} For any source or target $v_t$, TRAIN phase means and departures are
\begin{equation}
 \bar v_\phi=\frac{1}{n_{\mathrm{tr},\phi}}
 \sum_{t\in\mathrm{TRAIN}:\phi(t)=\phi}v_t,
 \qquad \delta v_t=v_t-\bar v_{\phi(t)} .
 \label{eq:supp-phase-centre}
\end{equation}
Stack $\delta x_t^\top$ and $\delta y_t^\top$ as rows of $X$ and $Z$. The road-coordinate ridge response is
\begin{equation}
 K_R^\star=\argmin_K\{\norm{Z-XK^\top}_F^2+\lambda\norm{K}_F^2\}
 =Z^\top X(X^\top X+\lambda I)^{-1},
 \label{eq:supp-ridge}
\end{equation}
with $\lambda=\kappa s_{\max}(XX^\top)$ and
$\kappa\in\{10^{-10},10^{-8},10^{-6},10^{-5},10^{-4},10^{-3},10^{-2},10^{-1},1,10,100\}$.
Validation selects $\kappa$ in the scoring frame of each comparison arm. The structural-characterization estimators use $\kappa=0.01$ in Porto, $0.1$ in San Francisco, $0.1$ in Chengdu's D frame, and $0.001$ for Chengdu's P-frame sensitivity. Each city is fitted separately; these settings are not a cross-city transfer rule.

\paragraph{P, D and CLIM scores.} Let $T_t=\UR^\top\Fcirc_t$ on ROAD support and $T_t=\Fcirc_t$ on FACE support. ROAD scoring projects full-face predictions into road coordinates; FACE scoring lifts road predictions to the complete routing field, so their denominators differ. Both fitted maps use TRAIN departures. P selects on raw VAL inputs and targets and predicts $B_{\mathrm P}x_t$ without a fitted intercept; D centres VAL/TEST with TRAIN phase means and predicts $B_{\mathrm D}\delta x_t$:
\begin{equation}
\begin{aligned}
 R_{\mathrm P}^2&=1-\frac{\sum_{t\in\mathrm{TEST}}\norm{T_t-B_{\mathrm P}x_t}^2}{\sum_{t\in\mathrm{TEST}}\norm{T_t}^2},\\
 R_{\mathrm D}^2&=1-\frac{\sum_{t\in\mathrm{TEST}}\norm{\delta T_t-B_{\mathrm D}\delta x_t}^2}{\sum_{t\in\mathrm{TEST}}\norm{\delta T_t}^2}.
\end{aligned}
\label{eq:supp-pd-scores}
\end{equation}
Numerators and denominators are pooled over TEST observations and folds within each city. The denominator is target energy, not variance about a TEST mean. CLIM is the source-free prediction $\bar T_{\phi(t)}$ on the same support and samples. For the P-frame physical display in Figure~\ref{fig:supp-6}, the observed field is $\UR T_t$ and the prediction is $\UR B_{\mathrm P}x_t$, using raw TEST inputs and targets without adding a phase climatology or substituting a D departure.

Full-face ridge fits $\Fcirc_t$ and projects its predictions with $P_0$. POD regression estimates its output basis from TRAIN departures only and selects rank and regularization on VAL in the common FACE metric. No TEST observation selects a basis, rank or regularization. These linear controls and the matched-capacity protocol in Note~\ref{note:spatial} are distinct experiments.

\paragraph{Nulls and uncertainty.} The rank-matched structural null fixes the ambient space and basis rank but not the subspace. Let $K_D$ be an orthonormal basis of $\ker D$. Each draw uses $G_{ij}\overset{\mathrm{iid}}{\sim}N(0,1)$, $Q=\operatorname{orth}(G)$ and $U_b=K_DQ$ at $\rank(\UR)$. The statistic is
\begin{equation}
 E(U)=1-\frac{\sum_t\norm{U^\top\Fcirc_t}^2}{\sum_t\norm{\Fcirc_t}^2},
 \qquad \Delta E_b=E(U_b)-E(\UR),
 \label{eq:supp-rank-null}
\end{equation}
evaluated on each city's full aligned observation pool without a source-response fit. It tests representation specificity, not source-based prediction.

The primary three-city D scores in Table~\ref{tab:crosscity} have paired-date bootstrap intervals and source--routing pairing-null results, reported in Note~\ref{note:coupling-results}. Dates are resampled with replacement within each fold; faces are not resampled. The pairing null permutes complete date blocks within TRAIN, VAL and TEST separately and refits the operator in every draw. Chengdu has only three dates, so its interval and null are descriptive. The neural paired-date contrast uses 9,999 bootstrap resamples, seed 20261003, dates resampled within folds and Holm adjustment within city/support. The matched-capacity implementation uses 10,000 multinomial paired-date resamples with seed 23 and summarizes the date-mean contrast. Chengdu's matched-capacity interval resamples only four alternate-hour TEST dates and is a small-sample descriptive interval.
